问题四与问题二、三采用相同的“日前线性规划--日内滚动执行--逐日审计”求解框架。新增的价格预测只提供各决策时点可用的价格路径，不改变 LP 的变量类型、约束处理方式或逐十分钟的执行逻辑。日前额度和多时点调整两条策略均先最小化经济目标；在货币容差内，再最小化电池吞吐量以消除等价循环。

\algorithmstart
\algostep{1}{构造因果信息}
在第 $d$ 日 $0{:}00$，仅读取此前已结束日的负荷、光伏、价格和 SOC 记录；以滚动 MAE 选择日初价格基准 $\widehat p^0$，并生成负荷、光伏与价格的联合残差池。
\algostep{2}{生成 Q4-2 日前额度}
对联合残差执行 K-medoids，保留 $8$ 条同日配对的代表路径；在 Q2 已冻结的 $(m,\alpha)$ 日历和风险储备下，求解式\eqref{eq:q4-2-dayahead}，锁定 $q_{d,1:T}$。
\algostep{3}{执行 Q4-2 日内 MPC}
每一时段仅读取当前实际 $(L,P,p)$ 与上期 SOC；根据已观测前缀更新后验权重和剩余预测，在锁定额度下求解 LP，并只实施当前一步。已锁定额度和已发生状态不回滚。
\algostep{4}{生成 Q4-3 日初承诺}
以 $\widehat L$、$\widehat P^0$ 和 $\widehat p^0$ 求解日初 LP，得到 $g^0_{d,t}$ 与储能预案；该承诺成为当天全部调整结算的唯一基准。
\algostep{5}{更新 Q4-3 未执行承诺}
到达 $6{:}00$、$12{:}00$、$18{:}00$ 时，用已发布预报和已观测价格前缀修正剩余预测；分别求解锁定原承诺与允许调整的 LP，仅当 $VoI_\tau$ 超过货币容差时替换未执行部分，并继续逐十分钟执行。
\algostep{6}{审计并导出结果}
核验信息截止点、联合残差配对、能量平衡、SOC、功率、额度、调整锁定和不同时充放电；随后按题设模板导出正式工作簿。每个 $\alpha$ 候选及 $K=4,8,12$ 均在同一闭环下复核实际成本、紧急购电量与耗时。
\algorithmend

主模型保持 LP；情景数固定为 $8$，日内 MPC 的视域为剩余时段加 $48$ 小时终端价值。价格维度加入后，仍须在 Q2 使用的求解时间预算内完成同一验证日的闭环评分，否则该候选不得进入正式部署。
